\section{Causal Recorder and Background Controls}
\label{sec:method}
\subsection{Reference policies}
\textbf{Uniform} retains the first available original observation in each of
$K=\lfloor(B-32)/36\rfloor$ equal time strata. Missing observations are not
interpolated. \textbf{Summaries} partition the epoch into
$\lfloor(B-32)/64\rfloor$ strata and record extrema and counts for nonempty
strata. They preserve bounds, not a raw trajectory. \textbf{Threshold} retains
observations whose causal score reaches or exceeds a fixed trigger boundary,
stopping when raw capacity is exhausted.

\subsection{Seeded detector-independent controls}
Reservoir sampling uses Algorithm R~\citep{vitter1985reservoir} to retain a uniform
sample without replacement of up to $K$ observations without knowing the final
count. Randomized-phase sampling draws integer $\phi$ uniformly in $[0,T)$ and
targets offsets $(jT+\phi)/K$. It takes the first available observation at or after
a target, skips targets crossed by gaps, and never duplicates an observation.
The latter is uniform in time, not equal-probability sampling of irregular
observations. Seeds 0--9 are all reported. A SHA-256 key per policy/seed/epoch
isolates skipped epochs; omitting the cap couples random draws across budgets
without requiring nested sample sets. An unused initial-fill reservoir draw
preserves that coupling without changing Algorithm R's marginal rule.

\subsection{Adaptive context with reserved background capacity}
The adaptive policy reserves $K_b=\lceil\rho K\rceil$ raw slots for background
observations and $K-K_b$ for triggered detail. Its baseline background rule takes
the first available observation in each of $K_b$ time strata independently of
anomaly scores. A trigger cannot consume this reservation, and unused slots are
not reclaimed by detail capture.

On a trigger, the recorder reserves room for the current observation and adds
the most recent affordable, not-yet-retained samples from its bounded pre-trigger
ring. A configured number of following samples may then be retained. Repeated
triggers can exhaust detail capacity but cannot exceed the encoded cap. Records
are deduplicated by timestamp and sorted at finalization. The ring and quotas
reset at epoch boundaries; finalized epochs cannot be reopened.

\subsection{Coverage and adjacent-pair controls}
The coverage control retains the same background quota but may replace its current
stratum's background candidate. It prefers a larger minimum squared distance from
previously retained background vectors, measured in frozen training-standardized
units. The first stratum supplies an initial anchor; records already requested as
trigger/context are protected from replacement. The control is a bounded diversity
heuristic, not a claimed novel or optimal sampling algorithm.

The pair control allocates $\lfloor K_b/2\rfloor$ background strata to adjacent
original observations. A pending sample is committed with its next native neighbor
only if both lie in the same stratum and epoch and their gap is at most 150 seconds.
Unpaired samples are not emitted as background, and an odd quota slot is left
unused. Trigger/context capture remains unchanged. At the 1,024-byte cap,
$K=27$ and $K_b=6$: singleton background spreads six samples across an hour,
whereas the pair control targets three consecutive pairs. This deliberately trades
temporal spread for inputs usable by a first-order forecasting receiver. It was
developed after inspecting the Mission1 temporal result.

\subsection{Causality and delivery}
An older buffered observation can be selected after a trigger because it was
already available and occupied bounded memory. Replacement also uses only the
observations seen so far. Packets are finalized at epoch end; capture-only
assessment therefore measures batching delay rather than immediate detection
reaction time. A separate finite-buffer simulator evaluates complete delivered
packets under declared contacts and loss. Its assumptions are distinguished from
the recording cap in Section~\ref{sec:delivery}.
